\problemname{Danska}

\noindent

Boschua ska flytta till Lund för att plugga på LTH. Du har upptäckt att Lund ligger i Skåne
och har därför blivit väldigt orolig över att Boschua ska börja låta dansk.

Du har hittat en gammal Dansk ordbok som har $M$ stycken ord i sig. Du vill nu visa för Boschua hur många
ord det finns som inte alls låter danska. Tyvärr tar det lång tid att räkna alla för hand, så du
behöver nu skriva ett program för detta.

Mer exakt, ett ord består av endast små engelska bokstäver \texttt{a}-\texttt{z}. Du vill då räkna antalet ord
av längd exakt $N$ som inte låter för danska. Ett ord låter för danskt ifall ordet innehåller minst en av orden
från ordboken i sig.

Exempelvis, anta att orboken endast består av ordet \texttt{knallert}. Då innehåller ordet \texttt{husknallertrumma}
ordet knallert, och anses därför vara för danskt.

Däremot innehåller ordet \texttt{knallerunt} \textbf{inte} ordet knallert på grund av förekomsten av ett \texttt{u},
och anses därför inte vara för danskt.

% Beräkna antalet strängar av längd $N$ som inte innehåller någon av de $M$ substrängarna modolu $10^9+7$.


\section*{Indata}
Den första raden i indatan innehåller heltalen $N$ och $M$ ($1 \le N \le 100$, $0 \le M \le 100$),
längden på orden vi betraktar och antalet ord i ordboken.

Därefter följer $M$ rader, där varje rad innehåller ett ord från ordboken. Alla dessa orden har längd som mest $N$
och består av bokstäver \texttt{a}-\texttt{z}. 


\section*{Utdata}
Skriv ut antalet ord av längd $N$ som inte låter för danskt. Eftersom detta talet kan bli väldigt stort,
skriv ut det \href{https://en.wikipedia.org/wiki/Modulo}{modulo} $10^9+7$. 

\section*{Poängsättning}
Din lösning kommer att testas på en mängd testfallsgrupper.
För att få poäng för en grupp så måste du klara alla testfall i gruppen.

\noindent
\begin{tabular}{| l | l | p{12cm} |}
  \hline
  \textbf{Grupp} & \textbf{Poäng} & \textbf{Gränser} \\ \hline
  $1$    & $10$       & $R, C \leq 30, Q=1$ \\ \hline
  $2$    & $13$       & $R, C \leq 30, Q \leq 10^4$ \\ \hline
  $3$    & $23$       & $Q=1$ \\ \hline
  $4$    & $32$       & Målet och lådans startposition är samma för alla frågor. \\ \hline
  $5$    & $32$       & Målet är samma för alla frågor. \\ \hline
  $6$    & $22$       & Inga ytterligare begränsningar. \\ \hline
\end{tabular}

\section*{Förklaring av exempelfall}
I första fallet finns det $26^3$ olika ord. Eftersom den Danska ordboken är tom i detta fallet förbjuds inga av dessa.
